\section{Current-state recurrences and the streaming bound}
\label{sec:streaming-consequences}

We now complete the streaming applications of the Gaussian determinant
test and the distance-bin test through two current-state iterations.
Both use one independent side matrix throughout the stream. They condition
on the present padded state, apply a complete bounded-prior block bound in
that conditional experiment, and only then average over state values.
Thus their information upper bounds never subtract two possibly infinite
averages of conditional total correlations.

\subsection{The Gaussian replica clock}

\begin{proposition}[Current-state iteration of Gaussian replicas]
\label{prop:gaussian-streaming}
Let $M(d)=o(d^2)$ be nonnegative and integer-valued, and let
$0<\epsilon(d)\le1/10$. A learner in
Definition~\ref{def:streaming-model} whose uniform-prior angular success is
at least $3/5$ satisfies
\[
 T(d)\ge c\,d\log(1/\epsilon(d))
\]
for an absolute $c>0$ and all sufficiently large $d$. The threshold may
depend on the memory sequence.
\end{proposition}

\begin{proof}
Use the first step of Lemma~\ref{lem:padding} to choose a fixed shared-seed
conditional experiment with success at least $3/5$ and then take its Borel
replacement. Its initialization is independent of the signal and all rows;
fresh transition and output randomness remains in the kernels. This
selection uses the original jointly measurable experiment and does not
require a jointly measurable family of Borel replacements.

Put $\beta=\log(1/\epsilon)$. If $T>d\beta$, the conclusion is immediate.
On the remaining branch, Lemmas~\ref{lem:padding} and
\ref{lem:early-clock} give padded states of entropy at most
\begin{equation}\label{eq:gaussian-clock-width}
 h_T=M\log2+\log(T+2)=o(d^2),
 \qquad \beta=o(d),\qquad T=o(d^2).
\end{equation}
Let
\[
 k=2\lfloor d/16\rfloor,\qquad t=k+1,\qquad
 \ell=4k,\qquad n=\lceil T/k\rceil.
\]
Append ignored fresh samples so that the padded stream has $n$ complete
blocks. Let $V_i$ be the padded state at the end of block $i$, with $V_0$
the data-independent initial padded state. The latter includes a possible
stopping decision at index zero. Draw a single standard Gaussian
$\ell\times d$ matrix $B$ independently of the entire selected learner
experiment and all appended rows, and put $Z=BS$.

Consider a positive-probability value $V_i=v$, and put
$p_v=\Prob\{V_i=v\}$. By Lemma~\ref{lem:conditional-priors}, its conditional
signal law has a Borel density $f_v$ satisfying
\[
 f_v\le L_v:=1/p_v
 \quad\hbox{relative to }\sigma_d,\qquad
 \sum_v p_v\log L_v=H(V_i).
\]
The next block matrix is independent of the pair $(S,V_i)$: the state uses
only the signal, completed rows, and learner randomness, and the new rows
are fresh. Given $v$, the next state is a finite kernel of the new matrix
$A$ and its exact label $AS$. The side matrix $B$ remains independent of
this entire conditional signal, block, and message experiment. Therefore
Theorem~\ref{thm:gaussian-block} applies with $L=L_v$ and with message
$V_{i+1}$.

Apply that numerical bound at each $v$ before averaging. Concavity of
$x\mapsto\log(2+x)$ and the entropy bound give
\begin{align}
 I(S;V_{i+1}\mid V_i,B,Z)
 &\le \frac{H(V_{i+1}\mid V_i)}{t}
      +Cd+C\sum_vp_v\log(2+\log L_v) \notag\\
 &\le \frac{h_T}{t}+Cd+C\log(2+H(V_i)) \notag\\
 &\le \frac{h_T}{t}+Cd+C\log(2+h_T).
 \label{eq:gaussian-current-step}
\end{align}
All conditional uses have bounded priors and finite conditional replica
correlation, as established in the block theorem. Since $t\asymp d$ and
$h_T=o(d^2)$, the last line is at most $C_1d$ for all sufficiently large
$d$, with one absolute $C_1$.

Write $\mathcal I(V)=I(S;V\mid B,Z)$. The chain rule and data processing
give
\[
 \mathcal I(V_{i+1})
 \le \mathcal I(V_i)+I(S;V_{i+1}\mid V_i,B,Z).
\]
Here $\mathcal I(V_0)=0$ by the independence of initialization and the
index-zero decision from $(S,B,Z)$. Iterating
\eqref{eq:gaussian-current-step} yields
\begin{equation}\label{eq:gaussian-information-clock}
 \mathcal I(V_n)\le C_1dn.
\end{equation}

The side row count satisfies $\ell=8\lfloor d/16\rfloor\le d/2$.
The postponed output is a kernel of $V_n$ with fresh randomness, so
\eqref{eq:endpoint-gaussian} gives
\begin{equation}\label{eq:gaussian-streaming-endpoint}
 \mathcal I(V_n)\ge
 \tfrac12(d-\ell-1)\log\frac1{5\epsilon}-\log2
 \ge c_1d\beta
\end{equation}
for large $d$. Lemma~\ref{lem:raw-sample-bound} also gives
$T>\lfloor d/2\rfloor$. In particular $T>0$, and for $d\ge32$ the
inequality $k\ge d/16$ gives the explicit clock comparison
\[
 dn\le d(T/k+1)\le16T+d\le18T.
\]
Combining this with \eqref{eq:gaussian-information-clock} and
\eqref{eq:gaussian-streaming-endpoint} proves the proposition.
\end{proof}

\subsection{The distance-bin clock}

\begin{proposition}[Current-state iteration of distance bins]
\label{prop:distance-streaming}
Under the hypotheses of Theorem~\ref{thm:streaming}, the distance-bin
block estimate \eqref{eq:distance-block} also gives
$T(d)\ge c\,d\log(1/\epsilon(d))$ for an absolute $c>0$ and all
sufficiently large $d$.
\end{proposition}

\begin{proof}
Apply Lemma~\ref{lem:padding} with $p=2/3$ and $p'=3/5$. It first fixes a
shared seed using the original jointly measurable experiment, then chooses
Borel versions in that conditional experiment and fixes its independent
initialization and remaining fresh randomness. We obtain deterministic
Borel state-based rules and an output with uniform-prior success at least
$3/5$. Pad stopping and postpone this output. The final output is a
function of the final tagged state.

For this route keep the convenient clock bound
\begin{equation}\label{eq:distance-clock-width}
 \mathfrak m_T=M\log2+\log(T+3).
\end{equation}
Indeed, the exact count in Lemma~\ref{lem:padding} is $(T+2)2^M$, so the
entropy of every padded state is at most $h_T\le\mathfrak m_T$.
Use
\[
 k=\lfloor d/16\rfloor,\qquad q=\lfloor d/8\rfloor,\qquad
 p=\lfloor d/2\rfloor,\qquad n=\lceil T/k\rceil,
\]
and append ignored fresh rows through the last block. Let $V_0,\ldots,V_n$
be the block-end states, starting with the data-independent tagged initial
state. Independently of the entire deterministic learner experiment and
all rows, draw one $p\times d$ standard Gaussian matrix $B$. It is used
only for analysis. Put $Z=BS$ and
$\mathcal D(V)=I(S;V\mid B,Z)$.

For a block from $V$ to $W$, condition on a value $v$ with probability
$p_v>0$. By Lemma~\ref{lem:conditional-priors}, the conditional signal density is bounded by
$\exp(b_v)$, where $b_v=\log(1/p_v)$. The next standard Gaussian matrix
$A$ is independent of $(S,V)$, the ending state is a function of
$(v,A,AS)$, and $B$ remains independent of the entire conditional
experiment. Apply Theorem~\ref{thm:distance-block} at this fixed $v$ and
then average. Since $\sum_vp_vb_v=H(V)$,
\begin{align}
 I(S;W\mid V,B,Z)
 &\le \frac{H(W\mid V)}q+C\left(d+\frac{H(V)}d\right) \notag\\
 &\le C_2\left(d+\frac{\mathfrak m_T}{d}\right).
 \label{eq:distance-current-step}
\end{align}
The last inequality uses $q\ge d/16$ for large $d$. In this calculation
the prior cost $b_v/d$ is averaged linearly; no entire state history is
retained. The chain rule gives
$\mathcal D(W)-\mathcal D(V)\le I(S;W\mid V,B,Z)$, and
$\mathcal D(V_0)=0$. Consequently
\begin{equation}\label{eq:distance-information-clock}
 \mathcal D(V_n)\le C_2 n\left(d+\frac{\mathfrak m_T}{d}\right).
\end{equation}

For the exact side row count $p=\lfloor d/2\rfloor$, the independent
residual-sphere endpoint \eqref{eq:endpoint-bins} gives
\begin{equation}\label{eq:distance-streaming-endpoint}
 \mathcal D(V_n)\ge
 \tfrac12(d-p-1)\log\frac1{4\epsilon}-\log2
 \ge c_2d\beta,\qquad \beta=\log(1/\epsilon),
\end{equation}
for sufficiently large $d$. This is the direct $90/d$ radius estimate
and the event of mass at least $1/2$ used in the distance-bin route.
The corresponding raw-data bound in Lemma~\ref{lem:raw-sample-bound}
also yields $T>\lfloor d/2\rfloor$, equivalently $T>d/2$ for integer $T$.
The zero-row experiment in that lemma has residual radius one.

If $T\ge d\beta$, the asserted lower bound is immediate. Otherwise
$\mathcal D(V_n)\le H(V_n)\le\mathfrak m_T$ and
\eqref{eq:distance-streaming-endpoint} give the route's clock inequality
\[
 c_2d\beta\le M\log2+\log(T+3)
 \le M\log2+\log(d\beta+3).
\]
Since $d\beta+3\le(d+3)(\beta+1)$ and
$\log(1+\beta)\le\beta$, this implies, for large $d$,
\begin{equation}\label{eq:distance-short-clock}
 \beta\le \frac{C(M+\log d+1)}d=o(d).
\end{equation}
Thus $T<d\beta=o(d^2)$ and $\mathfrak m_T=o(d^2)$. The right side of
\eqref{eq:distance-information-clock} is consequently at most $C_3nd$.
For $d\ge32$, $k\ge d/32$, and $T>d/2$ gives
\[
 dn\le32T+d\le34T.
\]
Comparing the resulting upper bound $\mathcal D(V_n)\le C_4T$ with
\eqref{eq:distance-streaming-endpoint} proves the proposition.
\end{proof}

\begin{proof}[Proof of Theorem~\ref{thm:streaming}]
Its uniform-prior success $2/3$ is at least the $3/5$ required by
Proposition~\ref{prop:gaussian-streaming}, which proves the theorem.
Proposition~\ref{prop:distance-streaming} gives a second proof using
the linear entropy cost. The later sections give three further proofs:
the Haar current-state argument in Proposition~\ref{haar:streaming}, and
the synthetic and incidence full-history arguments in
Proposition~\ref{syn:streaming} and Theorem~\ref{inc:precision}.
Each uses the same learner model, while
retaining its own side law, conditional experiment, and clock estimate.
\end{proof}
